\session{14}{11/21}{1限}{4}{最終発表：AI活用提案}

\goals{
  \item 選んだ企業・組織に対するAI活用提案を、課題と目的から始まる一貫した構成でまとめる。
  \item 提案作成でどうAIを使ったかを明示し、AIの寄与と自分の判断を分けて発表する。
  \item 相互評価を通じて、他者の提案の強みと改善点を見つけ、自分の提案に活かす。
}

\guidesec{解説}

\topic{AI活用提案の構成}
最終発表では、各自が選んだ企業・組織に対するAI活用提案を発表します。提案は、講義で学んだ内容を1つにまとめる場です。
次の構成を標準とします。
\par\noindent\begin{gtabh}{colspec={Q[l,wd=34mm] X[l] Q[l,wd=22mm]}}
  項目 & 内容 & 関連する回 \\
  1. 課題と目的 & 対象の企業・組織が抱える経営上の課題は何か。提案は何を達成するためのものか。「AIを使うこと」を目的にしない。 & 第1・7・14回 \\
  2. 現状 & 対象の業務プロセス、情報システム、データの現状。公開情報と、可能なら関係者への聞き取りに基づく。 & 第4・5・8回 \\
  3. 提案の内容 & どの業務の、どの作業に、どんなAIを、どう組み込むか。人が担う部分と人の関与の設計。 & 第5・6・11回 \\
  4. 実現可能性 & データの有無、コスト、組織の受け入れ、法規制の4観点。自社開発か外部サービスか。 & 第6・8・12・13回 \\
  5. 期待される効果と KPI & 導入前後で比較できる指標。効果の見込みの根拠。 & 第7・10回 \\
  6. リスクと対策 & バイアス、プライバシー、著作権、セキュリティ、誤情報。ガバナンスの体制。 & 第13回 \\
  7. 段階的な進め方 & PoC から本番展開までの段階と、各段階の判断基準。 & 第7・12回 \\
  8. AI活用記録 & 提案作成でAIをどう使ったか。AIの寄与と自分の判断・検証の区別。 & 第1回・付録 \\
\end{gtabh}

\topic{よい提案の条件}
\begin{itemize}
  \item \textbf{課題が具体的}　「業務の効率化」ではなく、「受注入力に1日3時間かかり、誤入力で月に10件の出荷ミスが起きている」のように、数値や事実で示す。
  \item \textbf{対象が絞られている}　全社的な構想より、1つの業務に絞った具体的な提案の方が、実現可能性と効果を検討できる。
  \item \textbf{人の役割が設計されている}　AIが何をして、人が何を確認し、誰が責任を持つかが明確。
  \item \textbf{根拠と推測が区別されている}　公開情報で確かめた事実と、自分の推測を分けて書く。推測には「推定」と明記する。
  \item \textbf{反対意見に答えている}　「なぜ今まで誰もやっていないのか」「失敗するとしたら何が原因か」に答えている。
\end{itemize}

\topic{発表でのAIの寄与と自分の判断の区別}
本講義の AI Policy のとおり、提案作成でAIをどう使ったかを明示します。隠す必要はなく、むしろ適切に使って検証した過程は評価の対象です。
発表では次のように区別して示します。
\begin{itemize}
  \item \textbf{AIの寄与}　アイデアの候補出し、公開情報の要約、構成案、文章の推敲、反対意見の生成など、AIが担った部分。
  \item \textbf{自分の判断}　課題の選定、提案の選択、評価基準、推奨の決定、AIの出力を採用・修正・不採用にした判断。
  \item \textbf{検証}　事実・数値・固有名詞を何で確かめたか。AIの出力の誤りをどう見つけ、どう直したか。
\end{itemize}
「AIの案をそのまま使った」部分があるなら、そう書いてかまいません。ただし、その部分の内容について質問されたら自分の言葉で説明できる必要があります。

\topic{相互評価の観点}
相互評価（ピアレビュー）は、順位を付けるためではなく、他者の視点から提案の強みと改善点を見つけ、提案をよりよくするために行います。
評価する側にとっても、提案を評価する観点を身につける学習の機会です。
\begin{checks}
  \item 課題と目的が明確で、AIの利用が目的化していないか。
  \item 提案の内容が具体的で、人の役割と責任が設計されているか。
  \item 実現可能性が4観点（データ、コスト、組織、法規制）で検討されているか。
  \item 効果が KPI で示され、根拠があるか。
  \item リスクと対策、段階的な進め方があるか。
  \item AIの寄与と自分の判断が区別され、検証の過程が示されているか。
  \item 最も良い点を1つ、改善すれば提案が強くなる点を1つ、具体的に指摘する。
\end{checks}

\topic{キーワード}
\par\noindent\begin{keywords}{}
  課題と目的 & 提案の出発点。AIを使うこと自体を目的にしない。 \\
  実現可能性の4観点 & データの有無、コスト、組織の受け入れ、法規制。 \\
  KPI & 導入前後で比較できる効果の指標。提案の説得力を決める。 \\
  AIの寄与と自分の判断の区別 & 発表での明示の仕方。AI Policy の「区別する」。 \\
  相互評価 & 他者の視点で強みと改善点を見つける評価。順位付けや批判が目的ではない。 \\
\end{keywords}

\guidesec{演習課題}

\exercise{1}{AI演習：提案の骨子を作り、AIの寄与と自分の判断を分けて記録する}
\instr{自分の提案の対象と課題を決めたうえで、次のプロンプトでAIと共同で骨子を作ってください。AIの出力に対して、必ず「この提案が失敗するとしたら何が原因か」「なぜ今まで誰もやっていないのか」を尋ね、答えを骨子に反映します。最後に、各項目について「AIの寄与」「自分の判断」「検証したこと」を表にします。}
\begin{promptbox}[提案の骨子作成]
あなたは経営コンサルタントです。私は大学の授業で「（企業・組織名）」へのAI活用提案を作っています。対象の課題は「（課題を具体的に）」です。提案の骨子を、(1) 課題と目的、(2) 現状、(3) 提案の内容（AIが担う部分と人が担う部分を分けて）、(4) 実現可能性（データ・コスト・組織・法規制）、(5) 効果と KPI、(6) リスクと対策、(7) 段階的な進め方、の7項目で、各項目3行以内で書いてください。私が示した事実と、あなたの推測を分けて書き、推測には「推定」と付けてください。
\end{promptbox}

\exercise{2}{ピアレビューを書く}
\instr{次の提案の要旨を読み、解説の相互評価の観点に沿って、「最も良い点」を1つ、「改善すれば提案が強くなる点」を2つ、具体的に書いてください。}
\begin{point}{提案の要旨：地域の書店チェーン（5店舗）への生成AI導入}
課題：店舗ごとの選書が店長の経験に依存し、売れ残りが多い。提案：全店に生成AIを導入し、選書・発注・POP作成・SNS 投稿を自動化する。
効果：業務効率化により売上向上。コスト：月額の利用料のみ。リスク：特になし。進め方：来月から全店で一斉に開始。
AI活用記録：ChatGPT で全文を作成。
\end{point}

\exercise{3}{発表の想定問答}
\instr{自分の提案について、発表で聞かれそうな厳しい質問を5つ挙げ、それぞれに答えを用意してください。質問の候補はAIに出させてかまいませんが、答えは自分で書きます。}

\guidesec{模範解答}

\begin{answer}{演習1}
骨子の例（地域の食品スーパー3店舗、日配品の廃棄削減）：
\par\noindent\begin{gtabh}{colspec={Q[l,wd=26mm] X[l] X[l,wd=26mm] X[l,wd=30mm]}}
  項目 & 内容（要約） & AIの寄与 & 自分の判断・検証 \\
  課題と目的 & 日配品の廃棄率が売上の3\%（業界平均の約1.5倍、推定）。廃棄率を2\%に下げる。 & 業界平均の目安を提示 & 廃棄率は店舗への聞き取り。業界平均は公開統計で確認し、「推定」と明記。 \\
  現状 & 発注は各店の担当者が経験で決定。POS はあるが予測には未使用。 & 現状の整理の枠組み & 聞き取りの内容を自分で整理。 \\
  提案の内容 & POS ベンダーの需要予測機能（SaaS）で発注の参考値を提示。担当者が天候・イベントを加味して確定。 & 自社開発案を提示したが不採用 & 規模から SaaS を選択。人の確定を残す設計は自分の判断。 \\
  実現可能性 & データ：POS 2年分あり。コスト：月額の追加料金。組織：担当者の経験を否定しない運用。法規制：個人情報は使わない。 & 4観点の整理 & 月額料金はベンダーのサイトで確認。 \\
  効果と KPI & 廃棄率、欠品率、担当者の発注時間。 & KPI の候補出し & 欠品率を加えたのは自分（廃棄だけ減らすと欠品が増えるため）。 \\
  リスクと対策 & 予測が外れる（人の確定で対応）。担当者の反発（参考値として段階導入）。 & リスクの列挙 & 「なぜ今までやっていないか」への答え（費用対効果が不明だった）を追加。 \\
  段階的な進め方 & 1店舗・日配品10品目で3か月試行。廃棄率が20\%以上下がれば全店へ。 & 段階の案 & 判断基準の数値は自分で設定。 \\
\end{gtabh}
\end{answer}

\begin{answer}{演習2}
\begin{itemize}
  \item 最も良い点：課題が「選書が店長の経験に依存し、売れ残りが多い」と具体的で、対象の業務（選書・発注）が明確である。
  \item 改善点1：提案の範囲が広すぎる（選書・発注・POP・SNS を一度に自動化）。売れ残りの削減という目的に直結する「発注量の参考値の提示」に絞り、店長の経験を活かす人の関与を設計すると、実現可能性と効果の検討ができる。
  \item 改善点2：効果・コスト・リスクの検討が不十分。「売上向上」ではなく、売れ残り率などの KPI と現状の数値を示す。コストには発注担当者の確認の工数と教育を含める。リスクには、予測の外れ、店長の反発、顧客データの扱いを挙げ、全店一斉ではなく1店舗での試行から始める。
        AI活用記録については、「ChatGPT で全文を作成」では自分の判断と検証が示されておらず、AI Policy の「区別する」を満たしていない。
\end{itemize}
\end{answer}

\begin{answer}{演習3}
例（食品スーパーの提案に対して）：
\begin{enumerate}[label=\arabic*.,leftmargin=2em]
  \item 「廃棄率3\%という数値の根拠は？」答え：店舗の担当者への聞き取りで、1店舗の直近3か月の実績。他の2店舗は未確認で、試行の際に全店で測定する。
  \item 「なぜ今まで導入していないのか？」答え：POS ベンダーの予測機能の存在を担当者が知らなかったことと、費用対効果が不明だったため。本提案は試行で効果を測ることを含む。
  \item 「担当者が参考値を無視したら？」答え：参考値と実際の発注の差を記録し、差が大きい場合の理由を月1回振り返る。参考値を強制しないのは、担当者の知識（地域の行事など）を活かすため。
  \item 「廃棄が減っても欠品が増えるのでは？」答え：欠品率も KPI に含め、両方を見て判断する。
  \item 「AIはどこで使ったのか？」答え：骨子の構成、リスクの列挙、業界平均の目安の提示。業界平均は公開統計で確認した。提案の対象の選定、SaaS の選択、KPI への欠品率の追加、判断基準の数値は自分の判断。
\end{enumerate}
\end{answer}

\guidesec{出席確認クイズの解説}
講義サイトの出席確認ページ（第14回）の5問です。
% 出席確認クイズ 第14回　最終発表：AI活用提案
%   attendance/attendance14.html から生成（解説は本ガイド用に加筆）。

\quizq{1}{AI活用提案で最初に明確にすべきことはどれでしょうか?}%
  {使うAIツールの名前}%
  {\correct{解決したい課題と目的}}%
  {予算だけ}%
  {発表スライドの色}%
  {正解は (b)。提案は「何のために」から始めます。解決したい課題と目的が明確でないと、AIを使うこと自体が目的化します（第7回の失敗要因）。(a) ツールの選定は目的が決まった後、(c) 予算と (d) スライドの色は本質ではありません。}

\quizq{2}{提案の実現可能性を検討する際の観点として適切なものはどれでしょうか?}%
  {見た目の派手さ}%
  {流行しているかどうかだけ}%
  {\correct{データの有無、コスト、組織の受け入れ、法規制}}%
  {発表時間の長さ}%
  {正解は (c)。実現可能性は、必要なデータの有無、コスト、組織の受け入れ、法規制などの観点から検討します。(a)(b)(d) は実現可能性と無関係です。この4観点は第8・12・6・13回で扱った内容に対応しています。}

\quizq{3}{提案の効果を示す方法として適切なものはどれでしょうか?}%
  {感想だけを書く}%
  {効果は示さなくてよい}%
  {\correct{導入前後で比較できる指標(KPI)を設定する}}%
  {導入費用の高さで示す}%
  {正解は (c)。効果は、導入前後で比較できる具体的な指標（KPI）を設定して示します（第7回）。(a) 感想や (d) 費用の高さでは効果を示せず、(b) 効果を示さない提案は採用されません。}

\quizq{4}{AIを使って作成した提案を発表するときの姿勢として適切なものはどれでしょうか?}%
  {自分の判断は不要である}%
  {すべてAIの出力のまま提出する}%
  {\correct{AIの寄与と自分の判断・検証を区別して示す}}%
  {AIを使ったことは隠す}%
  {正解は (c)。本講義の AI Policy のとおり、AI がどこに寄与し、自分が何を判断・検証したかを区別して示します。(a)(b) は自分の判断の放棄、(d) は透明性の否定です。}

\quizq{5}{相互評価(ピアレビュー)の目的として適切なものはどれでしょうか?}%
  {順位を付けることだけ}%
  {\correct{他者の視点で提案の強みと改善点を見つける}}%
  {形式的に済ませること}%
  {発表者を批判すること}%
  {正解は (b)。相互評価は、他者の視点から強みと改善点を見つけ、提案をよりよくするために行います。(a) 順位付けや (d) 批判が目的ではなく、(c) 形式的に済ませては意味がありません。評価する側も、提案を評価する観点を身につける学習機会になります。}

